\begin{textsample}{POS Dim 3 – rr – Score 16.00 – rr/rr\_inf\_5.txt}  \label{ex:f3_pos_225}
3.2 AS COMPETÊNCIAS E OS CAMPOS DE EXPERIÊNCIAS À LUZ DA BNCC O currículo territorial da Educação \textbf{Infantil} reitera a concepção que vincula educar e \textbf{cuidar}, entendendo o \textbf{cuidado} como algo indissociável do processo educativo.
Nesse contexto, as \textbf{creches} e \textbf{pré-escolas}, ao \textbf{acolher} as vivências e os conhecimentos construídos pelas \textbf{crianças} no ambiente familiar e no contexto de sua comunidade, e articulá -los em sua proposta pedagógica, têm o objetivo de ampliar o universo de experiências, conhecimentos e habilidades dessas \textbf{crianças}, valorizando as especificidades da etapa e reiterando as DCNEI.
O currículo do Estado de Roraima propõe que uma organização curricular para a Educação \textbf{Infantil} deve ser constituída a partir de cinco aspectos : - Princípios da Educação \textbf{Infantil} – Os três grandes princípios que devem guiar o projeto pedagógico da unidade de Educação \textbf{Infantil} : éticos ( autonomia, responsabilidade, solidariedade, respeito ao bem-comum, ao meio ambiente e às diferentes culturas, identidades e singularidades ) ; políticos ( direitos de cidadania, exercício da criticidade, respeito à ordem democrática ) ; estéticos ( sensibilidade, criatividade, ludicidade, liberdade de expressão nas diferentes manifestações artísticas e culturais ) ( DCNEI, Art. 6º ). - \textbf{Cuidar} e Educar – A indissociabilidade do educar e \textbf{cuidar}, pressuposto da Educação Básica como um todo, é um compromisso com a integralidade da educação dos sujeitos e de sensibilidade e responsabilidade com o futuro da humanidade e do planeta ( DCNEI, Art. 8º ). - Interações e \textbf{Brincadeiras} – Tendo em vista a \textbf{centralidade} do \textbf{brincar} e dos relacionamentos na vida das \textbf{crianças} \textbf{pequenas}, esses dois eixos possibilitam as aprendizagens, o desenvolvimento e a socialização das \textbf{crianças} na Educação \textbf{Infantil} ( DCNEI, Art. 9º ). - Seleção de práticas, saberes e conhecimentos – A seleção de práticas sociais, saberes e conhecimentos significativos e contextualmente relevantes para as novas gerações é uma obrigação da escola em uma sociedade complexa.
As experiências que emergem da vida cotidiana dão origem aos conhecimentos a serem compartilhados e reelaborados.
As propostas curriculares, em seus discursos e na sua operacionalização, também constituem as subjetividades de \textbf{crianças} e dos \textbf{adultos}, pois a formação pessoal e social não está dissociada da formação do mundo físico, natural e social ( DCNEI, Art. 8º e 9º ). - \textbf{Centralidade} das \textbf{crianças} – A atitude de acolhimento das singularidades dos \textbf{bebês} e das \textbf{crianças} e a criação de espaço para a constituição de culturas \textbf{infantis} definem a \textbf{centralidade} da \textbf{criança}.
As diversidades culturais, sociais, \textbf{etárias}, étnico-raciais, econômicas e políticas de suas famílias e comunidades, presentes em sua vida, precisam ser contempladas nos projetos educacionais ( DCNEI, Art. 4º ).
Os aspectos anteriormente citados embasam as relações pedagógicas, os \textbf{cuidados}, assim como os temas, as metodologias e proposições que constituem o modo de gestão dos grupos de \textbf{crianças} e da \textbf{instituição}, além da programação dos ambientes no \textbf{dia} a \textbf{dia} da unidade de Educação \textbf{Infantil}, ou seja, o seu currículo ( BRASIL, 2016, p. 58 ).
As concepções que fomentam as competências e habilidades nos campos de experiências visam : - \textbf{Acolher} vivência e conhecimentos construídos pelas \textbf{crianças} no ambiente familiar e contexto de sua comunidade ; - Articular a vivência e conhecimentos em suas propostas pedagógicas ; - Ampliar o universo de experiências, conhecimentos e habilidades ; - Diversificar e consolidar novas aprendizagens por meio da socialização, autonomia e comunicação ; - Atuar de maneira complementar à educação familiar ( \textbf{Bebês}, \textbf{Crianças} bem \textbf{pequenas} e \textbf{pequenas} ) ; - Fomentar aprendizagens bem próximas aos contextos familiar e escolar.

% matched lemmas: acolher, adulto, bebê, brincadeira, brincar, centralidade, creche, criança, cuidado, cuidar, dia, etário, infantil, instituição, pequeno, pré-escola
\end{textsample}
